\documentclass[12pt, letterpaper]{article}
\usepackage{graphicx} % Required for inserting images
%\usepackage[T1]{fontenc}
\usepackage{hyperref}
\usepackage{tasks}
\usepackage{amsmath}

\title{Programy użytkowe \\ Lista 11}
\date{}
\usepackage{polski}
\begin{document}
\maketitle
\begin{enumerate}
\item Korzystając z \href{https://pl.wikipedia.org/wiki/Uk%C5%82ad_r%C3%B3wna%C5%84_liniowych#Posta%C4%87_i_zapis}{postaci macierzowej} rozwiąż następujące układy równań
\begin{tasks}[label=\alph*)](2)
\task 
$
\begin{cases}
3x+2y-z=1\\
2x-2y+4z=-2\\
-2x+y-z=0
\end{cases}
\quad$

\task 
$\begin{cases}
3x+2y-z=1\\
2x-2y+4z=-2
\end{cases}$
\end{tasks}
Uwaga: w drugim przypadku mamy nieonaczony układ równań. Należy więc skorzystać z funkcji \verb|rref| na rozszerzonej macierzy układu równań i samodzielnie odwikłać rozwiązania ze zwróconej postaci schodkowej.
\item Narysuj wykresy funkcji $\sin{x}$ oraz $\cos{x}$. Aby znalazły się na jednym rysunku, między wywołaniami funkcji \verb|plot| użyj komendy \verb|hold on|. Zwróć uwagę, że w Octave musisz ręcznie zdefiniować wektor argumentów oraz wartości.
\item Korzystając z pętli \verb|for| nanieś na wykres kolejne funkcje Bessela $J_{\alpha_1},...,J_{\alpha_n}$.
\end{enumerate}
\textbf{Zadanie domowe.} Wybierz dowolną inną \href{https://docs.octave.org/v4.4.0/Special-Functions.html}{funkcję specjalną} niż $J_\alpha$, wykonaj jej wyjkres, umieś go w pliku w \LaTeX{u} i dołącz wyrażenie ją opisujące.
\end{document}